% ----------------------------------------------------------------
% AMS-LaTeX Paper ************************************************
% **** -----------------------------------------------------------
\documentclass{amsart}
%\usepackage{txfonts}
\usepackage{graphicx}
\usepackage{enumitem}
\usepackage{amsmath}
\usepackage{amssymb}
% ----------------------------------------------------------------
\vfuzz2pt % Don't report over-full v-boxes if over-edge is small
\hfuzz2pt % Don't report over-full h-boxes if over-edge is small
% THEOREMS -------------------------------------------------------
\newtheorem{thm}{Theorem}[section]
\newtheorem{cor}[thm]{Corollary}
\newtheorem{lem}[thm]{Lemma}
\newtheorem{prop}[thm]{Proposition}
\theoremstyle{definition}
\newtheorem{defn}[thm]{Definition}
\theoremstyle{Exercise}
\newtheorem{ex}[thm]{Exercise}
\theoremstyle{remark}
\newtheorem{rem}[thm]{Remark}
\theoremstyle{rule}
\newtheorem{rul}[thm]{Rule}

\numberwithin{equation}{section}
% MATH -----------------------------------------------------------
\newcommand{\norm}[1]{\left\Vert#1\right\Vert}
\newcommand{\abs}[1]{\left\vert#1\right\vert}
\newcommand{\set}[1]{\left\{#1\right\}}
\newcommand{\Real}{\mathbb R}
\newcommand{\Z}{\mathbb Z}
\newcommand{\To}{\longrightarrow}
\newcommand{\BX}{\bB(X)}
\newcommand{\A}{\mathcal{A}}
% ----------------------------------------------------------------

% define some simple, commonly-used commands
\newcommand{\eps}{\varepsilon}
\newcommand{\dsum}{\displaystyle\sum}
\newcommand{\dint}{\displaystyle\int}

\newcommand{\pdr}[2]{\dfrac{\partial{#1}}{\partial{#2}}}
\newcommand{\pdrr}[2]{\dfrac{\partial^2{#1}}{\partial{#2}^2}}
\newcommand{\pdrt}[3]{\dfrac{\partial^2{#1}}{\partial{#2}{\partial{#3}}}}
\newcommand{\dr}[2]{\dfrac{d{#1}}{d{#2}}}
\newcommand{\aver}[1]{\langle {#1} \rangle}
\newcommand{\Baver}[1]{\Big\langle {#1} \Big\rangle}

\newcommand{\bzero}{\mathbf 0}
\newcommand{\bGamma}{\mbox{\boldmath{$\Gamma$}}}
\newcommand{\btheta}{\boldsymbol \theta}
\newcommand{\bchi}{\mbox{\boldmath{$\chi$}}}
\newcommand{\bnu}{\boldsymbol \nu}
\newcommand{\bmu}{\boldsymbol \mu}
\newcommand{\brho}{\mbox{\boldmath{$\rho$}}}
\newcommand{\bxi}{\boldsymbol \xi}
\newcommand{\bnabla}{\boldsymbol \nabla}
\newcommand{\bOm}{\boldsymbol \Omega}
\newcommand{\blambda}{\boldsymbol \lambda}
\newcommand{\bsigma}{\boldsymbol \sigma}

\newcommand{\bbR}{\mathbb R}
\newcommand{\bbC}{\mathbb C}
\newcommand{\bbQ}{\mathbb Q}
\newcommand{\bbN}{\mathbb N}
\newcommand{\bbZ}{\mathbb Z}

\newcommand{\ba}{\mathbf a} \newcommand{\bb}{\mathbf b}
\newcommand{\bc}{\mathbf c} \newcommand{\bd}{\mathbf d}
\newcommand{\be}{\mathbf e} \newcommand{\bff}{\mathbf f}
\newcommand{\bg}{\mathbf g} \newcommand{\bh}{\mathbf h}
\newcommand{\bi}{\mathbf i} \newcommand{\bj}{\mathbf j}
\newcommand{\bk}{\mathbf k} \newcommand{\bl}{\mathbf l}
\newcommand{\bm}{\mathbf m} \newcommand{\bn}{\mathbf n}
\newcommand{\bo}{\mathbf o} \newcommand{\bp}{\mathbf p}
\newcommand{\bq}{\mathbf q} \newcommand{\br}{\mathbf r}
\newcommand{\bs}{\mathbf s} \newcommand{\bt}{\mathbf t}
\newcommand{\bu}{\mathbf u} \newcommand{\bv}{\mathbf v}
\newcommand{\bw}{\mathbf w} \newcommand{\bx}{\mathbf x}
\newcommand{\by}{\mathbf y} \newcommand{\bz}{\mathbf z}
\newcommand{\bA}{\mathbf A} \newcommand{\bB}{\mathbf B}
\newcommand{\bC}{\mathbf C} \newcommand{\bD}{\mathbf D}
\newcommand{\bE}{\mathbf E} \newcommand{\bF}{\mathbf F}
\newcommand{\bG}{\mathbf G} \newcommand{\bH}{\mathbf H}
\newcommand{\bI}{\mathbf I} \newcommand{\bJ}{\mathbf J}
\newcommand{\bK}{\mathbf K} \newcommand{\bL}{\mathbf L}
\newcommand{\bM}{\mathbf M} \newcommand{\bN}{\mathbf N}
\newcommand{\bO}{\mathbf O} \newcommand{\bP}{\mathbf P}
\newcommand{\bQ}{\mathbf Q} \newcommand{\bR}{\mathbf R}
\newcommand{\bS}{\mathbf S} \newcommand{\bT}{\mathbf T}
\newcommand{\bU}{\mathbf U} \newcommand{\bV}{\mathbf V}
\newcommand{\bW}{\mathbf W} \newcommand{\bX}{\mathbf X}
\newcommand{\bY}{\mathbf Y} \newcommand{\bZ}{\mathbf Z}

\newcommand{\cA}{\mathcal A} \newcommand{\cB}{\mathcal B}
\newcommand{\cC}{\mathcal C} \newcommand{\cD}{\mathcal D}
\newcommand{\cE}{\mathcal E} \newcommand{\cF}{\mathcal F}
\newcommand{\cG}{\mathcal G} \newcommand{\cH}{\mathcal H}
\newcommand{\cI}{\mathcal I} \newcommand{\cJ}{\mathcal J}
\newcommand{\cK}{\mathcal K} \newcommand{\cL}{\mathcal L}
\newcommand{\cM}{\mathcal M} \newcommand{\cN}{\mathcal N}
\newcommand{\cO}{\mathcal O} \newcommand{\cP}{\mathcal P}
\newcommand{\cQ}{\mathcal Q} \newcommand{\cR}{\mathcal R}
\newcommand{\cS}{\mathcal S} \newcommand{\cT}{\mathcal T}
\newcommand{\cU}{\mathcal U} \newcommand{\cV}{\mathcal V}
\newcommand{\cW}{\mathcal W} \newcommand{\cX}{\mathcal X}
\newcommand{\cY}{\mathcal Y} \newcommand{\cZ}{\mathcal Z}

%%%%%%%%%%%%%%Start%%%%%%%%%%%%%Start%%%%%%%%%%%Start%%%%%%%%%%%%%%%Start%%%%%%%%%%%%%%%%%%%%%%%%%Start%%%%%%%%%%%%%%%%
%%%%%%%%%%%%%%Start%%%%%%%%%%%%%Start%%%%%%%%%%%Start%%%%%%%%%%%%%%%Start%%%%%%%%%%%%%%%%%%%%%%%%%Start%%%%%%%%%%%%%%%%
%%%%%%%%%%%%%%Start%%%%%%%%%%%%%Start%%%%%%%%%%%Start%%%%%%%%%%%%%%%Start%%%%%%%%%%%%%%%%%%%%%%%%%Start%%%%%%%%%%%%%%%%
\usepackage{fancyhdr}

\pagestyle{fancy}
\fancyhf{}
\rhead{}
\chead{\includegraphics[scale=.1]{snhu_logo.png}}


\begin{document}
\title{\sf Module Two Problem Set}

\begin{center}
\includegraphics[scale=.1]{snhu_logo.png}
\end{center}

%\thm{bbjh}
\maketitle
This document is proprietary to Southern New Hampshire University. It and the problems within may not be posted on any non-SNHU website.\\\\\\\\
\begin{center}
%Enter your name below this line:
Tyler Ellis
\end{center}


\begin{center}
\rule{\textwidth}{0.4pt}
\end{center}


\newpage


\newpage

\section*{}
\section*{}

Directions: Type your solutions into this document and be sure to show all steps for arriving at your solution. Just giving a final number may not receive full credit.\\

\section*{Problem 1}
\subsection*{Part 1}
{\bf Indicate whether the argument is valid or invalid. For valid arguments, prove that the argument is valid using a truth table. For invalid arguments, give truth values for the variables showing that the argument is not valid.}\\
 \begin{enumerate}

\item \[
\begin{array}{||c||}
\hline \hline
(p \land q) \to r\\
\\
\therefore (p \lor q) \to r\\
\hline \hline
\end{array}
\]\\\\
 %Enter your answer below this comment line.
\item Invalid. A counterexample:
If \( p = \text{true}, q = \text{true}, r = \text{false} \). Then:
\[
(p \land q) \to r = \text{true} \to \text{false} = \text{false}
\]
\[
(p \lor q) \to r = \text{true} \to \text{false} = \text{false}
\]
However, the truth of the premise does not guarantee the conclusion. Not a tautology, thus this is invalid.
 \\\\
   \end{enumerate}

\subsection*{Part 2}
{\bf Converse and inverse errors are typical forms of invalid arguments. Prove that each argument is invalid by giving truth values for the variables showing that the argument is invalid. You may find it easier to find the truth values by constructing a truth table.}\\
 \begin{enumerate}[label=(\alph*)]
\item Converse error
\[
\begin{array}{||c||}
\hline \hline
p \to q\\
q\\
\\
\therefore p\\
\hline \hline
\end{array}
\]\\\\
\item Converse Error:
If \( p = \text{false}, q = \text{true} \). Then:
\[
p \to q = \text{true}, q = \text{true}, \text{ but } p = \text{false}
\]
Invalid argument.
 \\\\
\item Inverse error
\[
\begin{array}{||c||}
\hline \hline
p \to q\\
\neg p\\
\\
\therefore \neg q\\
\hline \hline
\end{array}
\]\\\\
\item Inverse Error:
Let \( p = \text{false}, q = \text{true} \). Then:
\[
p \to q = \text{true}, \neg p = \text{true}, \neg q = \text{false}
\]
So, \(\neg q\) is false even though both premises are true. Invalid.
\\\\
\end{enumerate}

\subsection*{Part 3}
{\bf Which of the following arguments are invalid and which are valid? Prove your answer by replacing each proposition with a variable to obtain the form of the argument. Then prove that the form is valid or invalid.}\\
 \begin{enumerate}[label=(\alph*)]
  \item \[
\begin{array}{||c||}
\hline \hline
\text{The patient has high blood pressure or diabetes or both.}\\
\text {The patient has diabetes or high cholesterol or both.}\\
\\
\therefore \text {The patient has high blood pressure or high cholesterol.
}\\
\hline \hline
\end{array}
\]\\\\\
\item Invalid.

\[
p = \text{high blood pressure}, q = \text{diabetes}, r = \text{high cholesterol}
\]
For:
\[
(p \lor q), (q \lor r) \not\Rightarrow (p \lor r)
\]
Counter:
Let \( p = \text{false}, q = \text{true}, r = \text{false} \). Gives true, thus the conclusion is false.
\\\\\

 
    \end{enumerate}
 \newpage
%--------------------------------------------------------------------------------------------------

\section*{Problem 2}
\subsection*{Part 1}


 Which of the following arguments are valid? Explain your reasoning.\\
 \begin{enumerate}[label=(\alph*)]
\item I have a student in my class who is getting an $A$. Therefore, John, a student in my class, is getting an $A$. \\\\
%Enter your answer below this comment line.
\item Invalid. Just because one student is getting an A doesn’t mean John is.
\\\\
\item Every Girl Scout who sells at least 30 boxes of cookies will get a prize. Suzy, a Girl Scout, got a prize. Therefore, Suzy sold at least 30 boxes of cookies.\\\\
%Enter your answer below this comment line.\item Invalid. Suzy may have received a prize for a different reason.
\\\\
 \end{enumerate}

\subsection*{Part 2}
Determine whether each argument is valid. If the argument is valid, give a proof using the laws of logic. If the argument is invalid, give values for the predicates $P$ and $Q$ over the domain ${a,\; b}$ that demonstrate the argument is invalid.\\
 \begin{enumerate}[label=(\alph*)]
\item \[
\begin{array}{||c||}
\hline \hline
\exists x\, (P(x)\; \land \;Q(x) )\\
\\
\therefore \exists x\, Q(x)\; \land\; \exists x \,P(x) \\
\hline \hline
\end{array}
\]\\\\
\item Valid. From \( \exists x (P(x) \land Q(x)) \), we know there is at least one \( x \) such that both predicates are true. So that same \( x \) satisfies both \( \exists x\,P(x) \) and \( \exists x\,Q(x) \). Valid.
 \\\\

\item \[
\begin{array}{||c||}
\hline \hline
\forall x\, (P(x)\; \lor \;Q(x) )\\
\\
\therefore \forall x\, Q(x)\; \lor \; \forall x\, P(x) \\
\hline \hline
\end{array}
\]\\\\
\item Invalid. Consider:
\[
P(a) = \text{true}, Q(a) = \text{false}, P(b) = \text{false}, Q(b) = \text{true}
\]
Then:
\[
\forall x (P(x) \lor Q(x)) = \text{true}
\]
But:
\[
\forall x\, Q(x) = \text{false}, \forall x\, P(x) = \text{false}, \Rightarrow (\forall x Q(x) \lor \forall x P(x)) = \text{false}
\]
 \\\\
 \end{enumerate}
 \newpage
%--------------------------------------------------------------------------------------------------


\section*{Problem 3}

Prove the following using a direct proof. Your proof should be expressed in complete English sentences.
\\\\

If $a$, $b$, and $c$ are integers such that $b$ is a multiple of $a^3$ and $c$ is a multiple of $b^2$, then $c$ is a multiple of $a^6$.
\\\\
%Enter your answer below this comment line.
Given:
\[
b = a^3k,\quad c = b^2m = (a^3k)^2m = a^6k^2m
\]
Thus, \( c \) is a multiple of \( a^6 \). 
\\\\

\newpage
%--------------------------------------------------------------------------------------------------
\section*{Problem 4}
Prove the following using a direct proof:
\\

The sum of the squares of 4 consecutive integers is an even integer.
%Enter your answer below this comment line.
For integers \( n, n+1, n+2, n+3 \). Then:
\[
n^2 + (n+1)^2 + (n+2)^2 + (n+3)^2 = 4n^2 + 12n + 14
\]
Even since all terms are even.
\\\\


 \newpage
 %--------------------------------------------------------------------------------------------------
\section*{Problem 5}

Prove the following using a proof by contrapositive:
\\\\

Let $x$ be a rational number. Prove that if $xy$ is irrational, then y is irrational.\\\\
%Enter your answer below this comment line.
Contrapositive: If \( y \) is rational, then \( xy \) is rational.

If \( x = \frac{a}{b}, y = \frac{c}{d} \), then:
\[
xy = \frac{ac}{bd} \in \mathbb{Q}
\]
Thus, if \( xy \) is irrational, then \( y \) must be irrational.
\\\\




 \newpage
%--------------------------------------------------------------------------------------------------

\section*{Problem 6}
Prove the following using a proof by contradiction:
\\\\


The average of four real numbers is greater than or equal to at least one of the numbers.
%Enter your answer below this comment line.
Assume average is less than all four numbers:
\[
\frac{a+b+c+d}{4} < a, b, c, d
\Rightarrow a+b+c+d < 4a, 4b, 4c, 4d
\]
Contradiction arises since sum can’t be less than all four values multiplied by 4. Therefore, at least one must be less than or equal to average.

 \newpage
%--------------------------------------------------------------------------------------------------

\section*{Problem 7}

Let $\displaystyle q = \frac{a}{b}$ and $\displaystyle r = \frac{c}{d}$ be two rational numbers written in lowest terms. Let $s = q + r$ and $\displaystyle s = \frac{e}{f}$ be written in lowest terms. Assume that $s$ is not $0$.\\

 Prove or disprove the following two statements.
\\\\
a.  If $b$ and $d$ are odd, then $f$ is odd.
\\\\
b. If $b$ and $d$ are even, then $f$ is even.
\\\\
%Enter your answer below this comment line.
a. True. Odd denominators added result in an odd denominator in lowest terms due to no factor of 2 to reduce.

b. False. Example:
\[
\frac{2}{4} + \frac{2}{4} = \frac{4}{4} = 1 = \frac{1}{1}
\]
Denominator is 1 (odd) even though both inputs were even.


\newpage
%--------------------------------------------------------------------------------------------------

\section*{Problem 8}
{\bf Define $P(n)$ to be the assertion that:}\\
\[\displaystyle \sum_{j=1}^{n}\, j^2 \;=\;\frac{n(n+1)(2n+1)}{6}\]\\\\
\begin{enumerate}[label=(\alph*)]
  \item Verify that $P(3)$ is true.\\\\
   %Enter your answer here.
   \[ 1^2 + 2^2 + 3^2 = 1 + 4 + 9 = 14, \quad \frac{3(3+1)(2\cdot3+1)}{6} = \frac{3\cdot4\cdot7}{6} = 14 \]
   \\\\

  \item Express $P(k)$.\\\\
   %Enter your answer here.
   \( P(k): \sum_{j=1}^k j^2 = \frac{k(k+1)(2k+1)}{6} \)
   \\\\

  \item Express $P(k+1)$.\\\\
   %Enter your answer here.
   \( P(k+1): \sum_{j=1}^{k+1} j^2 = \frac{(k+1)(k+2)(2k+3)}{6} \)
   \\\\

   \item In an inductive proof that for every positive integer $n$,
   \[\displaystyle \sum_{j=1}^{n}\, j^2 \;=\;\frac{n(n+1)(2n+1)}{6}\]
   what must be proven in the base case?\\\\
    %Enter your answer here.
    Base case: Show \( P(1): 1^2 = \frac{1(1+1)(2\cdot1+1)}{6} = 1 \)
    \\\\

    \item In an inductive proof that for every positive integer $n$,
   \[\displaystyle \sum_{j=1}^{n}\, j^2 \;=\;\frac{n(n+1)(2n+1)}{6}\]
   what must be proven in the inductive step?\\\\
   %Enter your answer here.
   Inductive step: Prove \( P(k) \Rightarrow P(k+1) \)
   \\\\

   \item What would be the inductive hypothesis in the inductive step from your previous answer?\\\\
    %Enter your answer here.
    Assume \( \sum_{j=1}^k j^2 = \frac{k(k+1)(2k+1)}{6} \)
    \\\\

   \item Prove by induction that for any positive integer n,
   \[\displaystyle \sum_{j=1}^{n}\, j^2 \;=\;\frac{n(n+1)(2n+1)}{6}\] \\\\
    %Enter your answer here.
    Proof:
    \[ \sum_{j=1}^{k+1} j^2 = \sum_{j=1}^k j^2 + (k+1)^2 = \frac{k(k+1)(2k+1)}{6} + (k+1)^2 \]
    \[ = \frac{(k+1)[k(2k+1) + 6(k+1)]}{6} = \frac{(k+1)(k+2)(2k+3)}{6} \]

\end{enumerate}

\end{document}